\documentclass{article}
\usepackage[utf8]{inputenc}
\usepackage[T1]{fontenc}
\usepackage{geometry}
\geometry{
    a4paper,
    left=2.5cm,
    right=2.5cm,
    top=2.5cm,
    bottom=2.5cm
}
\usepackage{graphicx}
\usepackage{float}
\usepackage{hyperref}
\usepackage[english]{babel}
\usepackage{amsmath}
\usepackage{amssymb}
\usepackage{array}
\usepackage{booktabs}
\usepackage{caption}
\usepackage{subcaption}
\usepackage{enumitem}
\usepackage{fancyhdr}
\usepackage{titlesec}
\usepackage{url}
\usepackage{textcomp}
\usepackage{xcolor}
\usepackage{listings}
\usepackage{multirow}
\usepackage{longtable}
\usepackage{tabularx}
\usepackage{times}

% Set up page style
\pagestyle{fancy}
\fancyhf{}
\rfoot{\thepage}
\renewcommand{\headrulewidth}{0pt}

% Title formatting
\titleformat{\chapter}[display]{\normalfont\bfseries}{\chaptertitle\thechapter}{0pt}{\Huge}
\titlespacing*{\chapter}{0pt}{50pt}{40pt}

% Document info
\title{Template Document}
\author{Academic Support Agent}
\date{\today}

\begin{document}
\maketitle

% Content from titlepage would be included here
 %% Titelseite zu bearbeiten in titlepage.tex


\newpage
\section{Grundlagen}
Referenzen auf Abbildungen werden mit Abbildung \ref{fig:fh_kiel} gebildet.\\
Referenzen auf Quellen werden mittels \cite{lit:datasheet} erstellt. Die Quellen sind in der "quellenangaben.bib" Datei hinzuzuf gen.\\
Empfohlen werden zur Benennung der Inhalte und leichterer Referenzierung folgende Formate:\\
Abbildungen = fig:abbildungsname\\
Tabellen = tab:tabellenname\\
Quellen = lit:literaturname\\
Anh nge = app:anhangsname\vspace*{3\baselineskip}
\lipsum[1] \cite{lit:datasheet}
%% Normale Grafiken
\begin{figure}[h!]
\centering
\includegraphics[width=0.5\textwidth]{Fachhochschule_Kiel}
\caption{\label{fig:fh_kiel}Beispielgrafik}
\end{figure}
\subsection{Grundlagen 1}
%% Textumflossene Grafiken
\par
\begin{wrapfigure}[4]{r}[0cm]{3cm}
\includegraphics[width=0.2\textwidth]{Fachhochschule_Kiel}
\label{fig:fh_kiel_floating}
\end{wrapfigure}
\lipsum[1] \cite{lit:schutzarten}
\subsection{Grundlagen 2}
\lipsum[1]


\newpage
\section{Durchf hrung}
\lipsum[1-2]
%% Tabelle
\begin{table}
\centering
\begin{tabular}{|l|l|l|}
   \hline
   Ger t & Typ  & Farbe\\
   \hline
   Fernseher & Sony & Mittelblau\\
   \hline
   \end{tabular}
 \caption{Bespieltabelle}
\end{table}


\newpage
\section{Auswertung}
\lipsum[1-5]
\newpage
\section{Fazit}
\lipsum[1-2]

% Content from appendix.tex would be included here
 %% Anh nge zu bearbeiten in appendix.tex
\end{document}